% Generated draft: requires mathematical and source review.
\documentclass[11pt]{article}
\usepackage[margin=25mm]{geometry}
\usepackage{fontspec}
\setmainfont{Latin Modern Roman}
\usepackage{amsmath,amssymb,mathtools,unicode-math}
\setmathfont{Latin Modern Math}
\usepackage{enumitem,needspace,xurl,hyperref}
\setlength{\parindent}{0pt}
\setlength{\parskip}{5pt}
\setlength{\emergencystretch}{3em}
\newcommand{\N}{\mathbb{N}}
\newcommand{\Z}{\mathbb{Z}}
\newcommand{\Q}{\mathbb{Q}}
\newcommand{\R}{\mathbb{R}}
\newcommand{\C}{\mathbb{C}}
\newcommand{\F}{\mathbb{F}}
\newcommand{\A}{\mathbb{A}}
\newcommand{\PP}{\mathbb{P}}
\newcommand{\E}{\mathbb{E}}
\newcommand{\eps}{\varepsilon}
\newcommand{\dd}{\,\mathrm{d}}
\begin{document}
\section*{UnsolvedMath ID: 30005419}
\textbf{Source problem code:} catalogue entry
\subsection{Definitions and mathematical statement}
For the module question, use the standard commutative local-algebra convention: \(R\) is a commutative Noetherian local ring with maximal ideal \(\mathfrak m\). Let \(D(R)\) be the derived category of complexes of \(R\)-modules. For \(X\in D(R)\), let \(\operatorname{thick}_R(X)\) denote the smallest full triangulated subcategory containing \(X\) and closed under direct summands. Interpret the source notation \(\langle X\rangle\) in this sense. A complex is perfect if it is isomorphic in \(D(R)\) to a bounded complex of finitely generated projective \(R\)-modules; these are precisely the objects of \(\operatorname{thick}_R(R)\).

Define ordinary support by
\[
\operatorname{Supp}_R X
 =\{\mathfrak p\in\operatorname{Spec}R:
 H^i(X)_{\mathfrak p}\ne 0\text{ for some }i\in\Z\}.
\]
Regard a finitely generated module \(M\) as a complex concentrated in degree zero. It is \emph{proxy small} if there exists a perfect complex \(P\) such that
\[
P\in\operatorname{thick}_R(M)
\quad\text{and}\quad
\operatorname{Supp}_R P=\operatorname{Supp}_R M.
\]
Equivalently, \(P\in\operatorname{thick}_R(M)\cap\operatorname{thick}_R(R)\) with the stated support equality. The interpretation of the angle brackets is a standard convention made explicit here; it is not defined in the supplied excerpt.

Call \(R\) a complete intersection if its \(\mathfrak m\)-adic completion has a presentation
\[
\widehat R\cong Q/(f_1,\ldots,f_c),
\]
where \(Q\) is a regular local ring and \(f_1,\ldots,f_c\) is a \(Q\)-regular sequence, with the empty sequence allowed.

\textbf{Main question.} Is the following implication valid for every commutative Noetherian local ring \(R\)?
\[
\begin{gathered}
\text{For every finitely generated }R\text{-module }M,
\text{ there exists }P\in
\operatorname{thick}_R(M)\cap\operatorname{thick}_R(R)\\
\text{with }\operatorname{Supp}_R P=\operatorname{Supp}_R M
\end{gathered}
\quad\Longrightarrow\quad
R\text{ is a complete intersection}.
\]
The supplied report excerpt attributes a characterization involving all \(M\in D^b(R)\) to Pollitz and asks whether it can be restricted to finitely generated modules. This attribution is retained without asserting an independently verified theorem.

\textbf{Further support-variety questions.} The original excerpt also asks when
\[
V_R(R)=\A^n
\]
and, more generally, how small \(V_R(M)\) can be. The registry formulation additionally asks which conical subvarieties occur as \(V_R(M)\) for \(M\in D^b(R)\). These questions are retained as qualified source targets: the supplied record does not define the cohomological support construction \(V_R\), its base field or ambient space, the integer \(n\), the intended measure of smallness, or whether \(D^b(R)\) requires finitely generated cohomology. Consequently, these further questions cannot be given a unique precise formulation from the supplied evidence alone.
\subsection{Short English statement}
Must a commutative Noetherian local ring be a complete intersection if every finitely generated module can build a perfect complex with the same ordinary support through finitely many triangulated operations and direct summands? The record also asks about the size and realizability of cohomological support varieties, but leaves essential conventions for those questions unspecified.
\subsection{Sources}
\begin{itemize}
\item Supplied local UnsolvedMath registry record, permanent ID \(30005419\), exact source problem code \texttt{catalogue entry}, title \emph{Proxy-Small Modules and Cohomological Support Varieties}; collection: Oberwolfach Reports Open Problems. The record supplies both a registry formulation and an original report excerpt.
\item \emph{Resolutions in Local Algebra and Singularity Theory} (2023), Oberwolfach Reports, pp.~327--395, DOI \texttt{10.4171/owr/2023/6}: \url{https://doi.org/10.4171/owr/2023/6}. The supplied excerpt carries the page marker 381 and attributes the bounded-complex characterization to Pollitz. No individual proposer is identified in the record.
\item Briggs--Grifo--Pollitz, \emph{Bounds on cohomological support varieties}: \url{https://arxiv.org/abs/2210.15574}. This reference is supplied as a literature-triage source; its results are not independently checked here.
\item Supplied literature assessment dated \texttt{2026-08-22}: the record labels the problem open and states that the restriction from bounded complexes to finitely generated modules remains unresolved, including after later work on bounds and realizability. This is supplied dated curation, not an independently verified literature assessment. The registry also labels its statement \texttt{corrected\_verified} and reports direct checking of the official report; those are supplied verification claims, not verification performed for this draft.
\end{itemize}
\end{document}
